\documentclass[aspectratio=169]{beamer}

% ---------- ENCODING & FONTS ----------
\usepackage[utf8]{inputenc}
\usepackage[T1]{fontenc}

% ---------- THEME & LOOK ----------
\usetheme{Madrid}
\usecolortheme{seahorse}
\usefonttheme{professionalfonts}
\setbeamertemplate{navigation symbols}{}
\setbeamercovered{transparent}
\setbeamertemplate{itemize items}[circle]
\setbeamertemplate{enumerate items}[default]
\setlength{\parskip}{0.25em}

% ---------- COLORS ----------
\usepackage{xcolor}
\definecolor{accent}{RGB}{0,92,184}
\definecolor{soft}{RGB}{233,244,255}
\setbeamercolor{structure}{fg=accent}
\setbeamercolor{title}{fg=white,bg=accent}
\setbeamercolor{frametitle}{fg=accent}
\setbeamercolor{block title}{bg=accent!12,fg=accent}
\setbeamercolor{block body}{bg=soft}

% ---------- PACKAGES ----------
\usepackage{booktabs}
\usepackage{amsmath, amssymb}
\usepackage{array}
\usepackage{multirow}
\usepackage{hyperref}
\hypersetup{colorlinks=true,linkcolor=accent,urlcolor=accent}

% ---------- SAFE FALLBACK FOR TRANSITIONS ----------
\makeatletter
\@ifundefined{transfade}{\newcommand{\transfade}[1][]{}}{}
\makeatother

% ---------- PLACEHOLDER BOX (compile-safe; no images required) ----------
\newcommand{\placeholderbox}[2][0.82\linewidth]{%
  \begin{center}
    \fbox{\begin{minipage}[c][0.46\textheight][c]{#1}
      \centering \vspace{0.25em}
      \textbf{#2}\\[0.6em]
      \emph{Insert figure/diagram here later.}
    \end{minipage}}
  \end{center}
}

% ---------- TITLE ----------
\title{\textbf{Optimal Compensation Systems, Deferred Compensation, and Mandatory Retirement \vspace{.5cm}\\  }}
\subtitle{ {\large ECON 300 - Labour Economics}}
%\institute{UNBC}
\author{Muhebullah Karimzada}

\date{Winter 2026}



\begin{document}

% ============================================================
% TITLE
% ============================================================
\begin{frame}[plain]
  \titlepage
\end{frame}

% ============================================================
\section{Learning Objectives \& Main Questions}
% ============================================================

\begin{frame}{Learning Objectives}
\transfade
\begin{itemize}[<+->]
  \item Discuss why seemingly inefficient compensation arrangements can be optimal for aligning employer--employee incentives.
  \item Explain why compensation arrangements may change across time/countries/workplaces.
  \item Outline the ``new economics of personnel'' and its insights for HR practices.
  \item Engage in debates around CEO pay and mandatory retirement.
\end{itemize}
\end{frame}

\begin{frame}{Main Questions (1/2)}
\transfade
\begin{itemize}[<+->]
  \item How do internal labour markets differ from competitive external markets?
  \item Why do wages rise with seniority? Does productivity rise continually?
  \item Why are superstars paid so much more than average performers?
  \item When are tournament-style pay structures helpful vs.\ harmful?
\end{itemize}
\end{frame}

\begin{frame}{Main Questions (2/2)}
\transfade
\begin{itemize}[<+->]
  \item What is mandatory retirement -- age discrimination or efficient institution?
  \item Why are CEO/executive salaries very high, and why have they risen?
  \item Why does academic tenure exist? When should firms promote from within vs.\ hire outside?
  \item When do organizations match outside offers -- and what about the winner's curse?
\end{itemize}
\end{frame}

% ============================================================
\section{Agency Theory \& Efficiency Wages}
% ============================================================

\begin{frame}{Agency Theory}
\transfade
\begin{itemize}[<+->]
  \item Principal--agent problem: design efficient contracts when incentives to cheat exist and information is asymmetric.
  \item Elicit ``truth-telling'' via explicit (collective agreements) or implicit contracts.
\end{itemize}
\end{frame}

\begin{frame}{Efficiency Wage Theory}
\transfade
\begin{itemize}[<+->]
  \item Wages affect productivity; paying above spot can raise effort, reduce shirking/turnover.
  \item At subsistence levels, higher wages can improve nutrition and productivity.
  \item Result: \alert{efficiency gains} from higher wages in certain settings.
\end{itemize}
\end{frame}

% ============================================================
\section{Superstars, Tournaments, and Teams}
% ============================================================

\begin{frame}{Economics of Superstars}
\transfade
\begin{itemize}[<+->]
  \item Small skill differences can be magnified when output is broadcast or scaled (concerts, TV shows, software platforms).
  \item Tiny per-person value improvements aggregate to large total value when audience or co-worker base is large.
\end{itemize}
\end{frame}

\begin{frame}{Salaries as Tournament Prizes}
\transfade
\begin{itemize}[<+->]
  \item Promotions act as \textbf{prizes}: large salary jumps for winners.
  \item Pros: stronger incentives, especially when few promotion chances remain.
  \item Cons: discourages cooperation, may foster sabotage or corruption if prize spreads are excessive.
\end{itemize}
\end{frame}

\begin{frame}{Efficient Pay Equality, Teams, and Cooperation}
\transfade
\begin{itemize}[<+->]
  \item Some inequality is necessary for incentives.
  \item \alert{Too much inequality} harms teamwork and invites sabotage.
  \item Solutions: tailor incentive schemes by level, base managers' bonuses on group output, create separate units, use external promotion margins, and HR policies that reward cooperation.
\end{itemize}
\end{frame}

\begin{frame}{The 1/N Problem (Team Averages)}
\transfade
\begin{itemize}[<+->]
  \item Paying by average team output creates \textbf{free-rider} issues.
  \item Large teams $\Rightarrow$ weak individual incentives; small teams $\Rightarrow$ peer pressure and better monitoring.
\end{itemize}
\end{frame}

\begin{frame}{Up-or-Out Rules}
\transfade
\begin{itemize}[<+->]
  \item Evaluations at set career points lead to promotion or exit (e.g., academic tenure).
  \item Can be inefficient if good-but-not-promoted workers leave; remedy: junior non-promoted slots to continue evaluation.
\end{itemize}
\end{frame}

\begin{frame}{Raiding, Offer-Matching, and the Winner's Curse}
\transfade
\begin{itemize}[<+->]
  \item Asymmetric information drives \textbf{raiding} for talent.
  \item Incumbent firms match if the worker is truly worth the outside offer; otherwise, they do not.
  \item Raiders risk \textbf{winner's curse} -- overpaying when information is noisy.
\end{itemize}
\end{frame}

\begin{frame}{Piece Rates: Positives and Negatives}
\transfade
\begin{itemize}[<+->]
  \item \textbf{Benefits:} rewards high output, attracts strong workers, saves monitoring.
  \item \textbf{Drawbacks:} output measurement can be hard/costly; weak when individuals have limited control over output.
\end{itemize}
\end{frame}

% ============================================================
\section{Deferred Compensation}
% ============================================================

\begin{frame}{Deferred Wages: Concept}
\transfade
\begin{itemize}[<+->]
  \item Wages rise with seniority even if productivity is flat.
  \item Raises and pension accruals backload pay -- requires long-term commitment.
  \item \textbf{Efficiency rationale:} boosts effort/honesty, reduces monitoring and unwanted turnover, screens out low-effort workers, lowers hiring/training costs.
\end{itemize}
\end{frame}

\begin{frame}{Figure 13.1(a) -- Wage Productivity Profiles: Deferred Wages Contract (Placeholder)}
\transfade
\placeholderbox{Figure 13.1(a) -- Deferred wages contract market.}
\end{frame}

\begin{frame}{Figure 13.1(b) -- Wage Productivity Profiles: Spot Market (Placeholder)}
\transfade
\placeholderbox{Figure 13.1(b) -- Spot market.}
\end{frame}

\begin{frame}{Figure 13.1(c) -- Wage Productivity Profiles: Deferred Wages with Company--Specific Training (Placeholder)}
\transfade
\placeholderbox{Figure 13.1(c) -- Deferred wages with company--specific training.}
\end{frame}

\begin{frame}{Figure 13.1(d) -- Wage Productivity Profiles: General Training (Placeholder)}
\transfade
\placeholderbox{Figure 13.1(d) -- General training.}
\end{frame}

\begin{frame}{WORKED EXAMPLE -- Deferred Compensation}
\transfade
\small
Assume a worker's VMP is \$10 per hour.\par
Deferred compensation system: \(W = 5 + 0.2T\), where \(T\) is tenure in years.\par
Worker starts at age 20; ignore discounting.\par
\begin{itemize}[<+->]
  \item \textbf{Break-even tenure:} set \(W=\text{VMP}\): \(5 + 0.2T = 10 \Rightarrow T=25\).
  \item \textbf{Implied mandatory retirement age:} \(45 + 25 = 70\).
\end{itemize}
\end{frame}

\begin{frame}{Underpayment or Overpayment (Exact Table)}
\transfade
\scriptsize
\begin{tabular}{lccccc}
\toprule
\textbf{Underpayment or Overpayment} & \textbf{Age} & \textbf{Tenure, T} & \textbf{VMP} & \textbf{W = 5 + 0.2T} & \textbf{W -- VMP} \\
\midrule
 & 20 & 0  & 10 & 5  & $-5$ \\
 & 30 & 10 & 10 & 7  & $-3$ \\
 & 40 & 20 & 10 & 9  & $-1$ \\
 & 50 & 30 & 10 & 11 & $+1$ \\
 & 60 & 40 & 10 & 13 & $+3$ \\
 & 70 & 50 & 10 & 15 & $+5$ \\
\bottomrule
\end{tabular}
\end{frame}

\begin{frame}{WORKED EXAMPLE -- Piece Rates and Team Compensation (Exact Table)}
\transfade
\scriptsize
\begin{tabular}{cccccc}
\toprule
\multicolumn{4}{c}{\textbf{Disutility of Effort}} & \multicolumn{2}{c}{\textbf{Pay (wage $\times$ effort)}} \\
\textbf{Unit of effort = output (1)} & \textbf{Wage per unit of effort (2)} & \textbf{You (3)} & \textbf{EI Slacker (4)} & \textbf{You (5)} & \textbf{EI Slacker (6)} \\
\midrule
1 & 10 & 5  & 7   & 10 & 10 \\
2 & 10 & 6  & 8   & 20 & 20 \\
3 & 10 & 7  & 9   & 30 & 30 \\
4 & 10 & 8  & 10  & 40 & 40 \\
5 & 10 & 9  & 11  & 50 & 50 \\
6 & 10 & 10 & 12  & 60 & 60 \\
7 & 10 & 11 & 13  & 70 & 70 \\
\bottomrule
\end{tabular}
\end{frame}

\begin{frame}{Do Workers Prefer Deferred Wages?}
\transfade
\begin{itemize}[<+->]
  \item Many workers prefer rising wage profiles as \textbf{forced saving}.
  \item Anticipating future consumption itself yields utility (consumption smoothing).
  \item Workers may willingly trade lower current wages for higher future wages.
\end{itemize}
\end{frame}

% ============================================================
\section{Executive Compensation}
% ============================================================

\begin{frame}{Executive Compensation: Stylized Facts}
\transfade
\begin{itemize}[<+->]
  \item Executive pay has risen relative to average worker pay.
  \item Typical packages combine base salary with stock options and equity.
  \item Board dynamics matter: directors may have incentives to award high pay (many are executives elsewhere).
\end{itemize}
\end{frame}

% ============================================================
\section{Mandatory Retirement}
% ============================================================

\begin{frame}{Rationales for Mandatory Retirement}
\transfade
\begin{itemize}[<+->]
  \item Enables deferred-wage contracts so PV(wages) = PV(productivity) to the firm.
  \item Provides finality to long-term arrangements and an \alert{efficient compensation rule}.
  \item Opens promotion and employment opportunities for younger workers.
  \item Increases planning certainty for staffing, pensions, and medical costs; aids workers' pre-retirement planning.
  \item Reduces monitoring costs for older workers where effort is hard to observe.
\end{itemize}
\end{frame}

\begin{frame}{Figure 13.2 -- Mandatory Retirement Age (Placeholder)}
\transfade
\placeholderbox{Figure 13.2 -- Mandatory Retirement Age.}
\end{frame}

\begin{frame}{Arguments Against Mandatory Retirement}
\transfade
\begin{itemize}[<+->]
  \item Human-rights concerns: potential age discrimination and efficiency loss if productive workers are forced out.
  \item May improve viability of pension systems, but trade-offs exist.
\end{itemize}
\end{frame}

\begin{frame}{Impact of Banning Mandatory Retirement}
\transfade
\begin{itemize}[<+->]
  \item Age--earnings profiles may \textbf{track productivity} more closely (higher wages for younger, lower for older).
  \item Deferred-wage systems may change; more frequent monitoring/evaluation.
  \item Possible reduction in training and fewer openings for younger workers.
  \item Forecasting/planning becomes more difficult; jobs may be redesigned to accommodate older workers.
\end{itemize}
\end{frame}

% ============================================================
\section{Chapter Summary}
% ============================================================

\begin{frame}{Summary}
\transfade
\begin{itemize}[<+->]
  \item Core ideas: agency and efficiency wages; superstars; tournaments; team incentives.
  \item Deferred compensation aligns incentives, screens workers, and supports long-term matches.
  \item Executive pay mixes salary and equity; governance shapes outcomes.
  \item Mandatory retirement has efficiency rationales but also equitable and efficiency costs; design matters.
\end{itemize}
\end{frame}

\begin{frame}
\centering
\Huge{\textbf{Thank You!}} \\
\end{frame}

\end{document}
